% !TEX program = xelatex
\PassOptionsToPackage{hidelinks}{hyperref}

\documentclass{resume}
\usepackage{zh_CN-Adobefonts_external} % Simplified Chinese Support using external fonts (./fonts/zh_CN-Adobe/)
% \usepackage{zh_CN-Adobefonts_internal} % Simplified Chinese Support using system fonts
\usepackage{linespacing_fix} % disable extra space before next section
\usepackage{cite}

\usepackage{setspace}
\setstretch{1.2}
\usepackage{fontawesome5}

%-------------------------
% Custom commands
\newcommand{\resumeItem}[2]{
  \item{
    \textbf{#1}{: #2}
  }
}
\newcommand{\resumeItemSimple}[1]{
  \item\small{
    {#1}\vspace{-2pt}
  }
}
\newcommand{\resumeItemSimpleNoBullet}[1]{
  \vspace{-5pt}\item[]\small{
    {#1}\vspace{-2pt}
  }
}

\newcommand{\Paper}[3]{
  \vspace{-1pt}\item \small
    \hangindent=1.4em \hangafter=1
    \textbf{#1}. #2
    \par{\footnotesize #3}
  \vspace{-4pt}
}

\newcommand{\PaperTLDR}[4]{
  \vspace{-1pt}\item \small
    \hangindent=1.4em \hangafter=1
    \textbf{#1}. #2
    \par{\footnotesize #3}
    \if\relax\detokenize{#4}\relax\else
      \par{\scriptsize \emph{TLDR: #4}}
    \fi
  \vspace{-4pt}
}

\newcommand{\ExperienceListStart}{
    \vspace{-1pt}\begin{tabular*}{0.97\textwidth}{l@{\extracolsep{\fill}}r}
}

\newcommand{\ExperienceListEnd}{
    \end{tabular*}\vspace{-5pt}
}

\newcommand{\Experience}[2]{
    $\bullet$ {\small#1} & {\small #2} \\
}

\newcommand{\AwardListStart}{
    \vspace{-1pt}\begin{tabular*}{0.97\textwidth}{l@{\extracolsep{\fill}}r}
}

\newcommand{\AwardListEnd}{
    \end{tabular*}\vspace{-5pt}
}

\newcommand{\Award}[2]{
    $\bullet$ {\small#1} & \textit{\small #2} \\
}

% Just in case someone needs a heading that does not need to be in a list
\newcommand{\resumeHeading}[4]{
    \begin{tabular*}{0.99\textwidth}[t]{l@{\extracolsep{\fill}}r}
      \textbf{#1} & #2 \\
      \textit{\small#3} & \textit{\small #4} \\
    \end{tabular*}\vspace{-5pt}
}

\newcommand{\resumeSubheading}[4]{
  \vspace{-1pt}\item
    \begin{tabular*}{0.97\textwidth}[t]{l@{\extracolsep{\fill}}r}
      \textbf{#1} & #2 \\
      {#3} & {#4} \\
    \end{tabular*}\vspace{-5pt}
}

\newcommand{\resumeSubheadingSimple}[2]{
  \vspace{-1pt}\item
    \begin{tabular*}{1.00\textwidth}[t]{l@{\extracolsep{\fill}}r}
      \large{#1} & \large{#2} \\
    \end{tabular*}%\vspace{-7pt}

%
}

\newcommand{\resumeSubSubheading}[2]{
    \begin{tabular*}{0.97\textwidth}{l@{\extracolsep{\fill}}r}
      \textit{\small#1} & \textit{\small #2} \\
    \end{tabular*}\vspace{-5pt}
}

\newcommand{\resumeSubItem}[2]{\resumeItem{#1}{#2}}

\newcommand{\resumeSubItemSimple}[1]{\resumeItemsimple{#1}}

\renewcommand{\labelitemii}{$\circ$}

\newcommand{\resumeSubHeadingListStart}{\begin{itemize}[leftmargin=*]}
\newcommand{\resumeSubHeadingListEnd}{\end{itemize}}
\newcommand{\resumeItemListStart}{\begin{itemize}}
\newcommand{\resumeItemListEnd}{\end{itemize}\vspace{-5pt}}
%
\begin{document}
\pagenumbering{gobble} % suppress displaying page number

%----------HEADING-----------------
\begin{center}
{\LARGE \scshape \bfseries 马琰} \\ \vspace{2pt}
(+86) 152-6982-9096 $\bullet$
\href{mailto:yanma23@m.fudan.edu.cn}{yanma23@m.fudan.edu.cn} $\bullet$
\href{https://yanmaaaa.github.io}{主页} $\bullet$
\href{https://scholar.google.com/citations?user=tnJyAAoAAAAJ\&hl=en}{Google Scholar}
\end{center}

%-----------研究兴趣-----------------
\section{研究兴趣}
我的研究兴趣位于强化学习与多模态模型的交叉方向。2025 年，我在视觉语言模型的强化学习训练方面积累了较为系统的经验，重点关注单视觉任务、多视觉任务以及 Agentic 视觉工具使用场景下的 RL Scaling。此前，我也从事过统一多模态模型与文本生成相关研究。当前我主要关注多模态生成模型的预训练。

%-----------教育经历-----------------
\section{教育经历}
  \resumeSubHeadingListStart
    \resumeSubheading
      {复旦大学，计算机科学技术博士在读}{上海，中国}
      {导师：\href{https://scholar.google.co.jp/citations?user=oIz_CYEAAAAJ\&hl=en}{刘鹏飞} 教授、\href{https://mmlab.siat.ac.cn/yuqiao}{乔宇} 教授}{2023.09 -- 2027.06（预计）}
    \resumeSubheading
      {复旦大学，计算机应用技术硕士}{上海，中国}
      {导师：\href{https://scholar.google.co.uk/citations?user=AG3OXS0AAAAJ\&hl=en}{李伟} 教授}{2020.09 -- 2023.06}
    \resumeSubheading
      {大连理工大学，计算机科学与技术学士}{大连，中国}
      {GPA：4.076/5，排名：12/123}{2016.09 -- 2020.06}
  \resumeSubHeadingListEnd

%-----------实习经历-----------------
\section{实习经历}
  \ExperienceListStart
    \Experience{\textbf{MiniMax \& HailuoAI} - 实习生；参与 M 系列基础模型与 Hailuo 视频生成模型研发}{2025.01 -- 2026.02}
    \Experience{\textbf{Shanghai AI Laboratory} - 实习生；文本生成与统一多模态模型研究}{2023.07 -- 2025.01}
  \ExperienceListEnd

%-----------论文发表-----------------
\section{论文发表}
  \resumeSubHeadingListStart
  \PaperTLDR
  {What Does Vision Tool-Use Reinforcement Learning Really Learn? Disentangling Tool-Induced and Intrinsic Effects for Crop-and-Zoom \href{https://arxiv.org/abs/2602.01334}{[\faFilePdf]}\href{https://github.com/GAIR-NLP/Med}{[\faGithub]}}
  {\textit{arXiv:2602.01334}}
  {\underline{Yan Ma}$^\dagger$, Weiyu Zhang$^\dagger$, Tianle Li, Linge Du, Xuyang Shen, Pengfei Liu}
  {我们发现当前视觉工具使用 RL 的性能提升主要来自模型内在能力增强与工具引入误差的减少，而非真正掌握了基于工具的纠错能力。}
  \PaperTLDR
  {One RL to See Them All: Visual Triple Unified Reinforcement Learning \href{https://arxiv.org/abs/2505.18129}{[\faFilePdf]}\href{https://github.com/MiniMax-AI/One-RL-to-See-Them-All}{[\faGithub]}}
  {\textit{arXiv:2505.18129}}
  {\underline{Yan Ma}$^\dagger$, Linge Du$^\dagger$, Xuyang Shen$^\dagger$, Shaoxiang Chen, Pengfei Li, Qibing Ren, Lizhuang Ma, Yuchao Dai, Pengfei Liu, Junjie Yan}
  {我们构建了统一的 RL 系统（V-Triune），联合训练视觉推理与感知能力，并在两类任务上均取得了广泛提升。}
  \PaperTLDR
  {Rethinking RL Scaling for Vision Language Models: A Transparent, From-Scratch Framework and Comprehensive Evaluation Scheme \href{https://arxiv.org/abs/2504.02587}{[\faFilePdf]}\href{https://github.com/GAIR-NLP/MAYE}{[\faGithub]}}
  {\textit{arXiv:2504.02587}}
  {\underline{Yan Ma}, Steffi Chern, Xuyang Shen, Yiran Zhong, Pengfei Liu}
  {我们提出了一个透明、可复现的 VLM 强化学习训练与评测框架，并系统分析了 RL 相对 SFT 的泛化优势及关键训练动态。}
  \PaperTLDR
  {ANOLE: An Open, Autoregressive, Native Large Multimodal Models for Interleaved Image-Text Generation \href{https://arxiv.org/abs/2407.06135}{[\faFilePdf]}\href{https://iclr-blogposts.github.io/2025/blog/fine-tuning-token-based-large-multimodal-models/}{[\faGlobe]}\href{https://github.com/GAIR-NLP/Anole}{[\faGithub]}}
  {\textit{arXiv:2407.06135; extended version accepted to ICLR 2025 Blog Track}}
  {Ethan Chern*, Jiadi Su*, \underline{Yan Ma*}, Pengfei Liu}
  {}
  \PaperTLDR
  {MoPS: Modular Story Premise Synthesis for Open-Ended Automatic Story Generation \href{https://arxiv.org/abs/2406.05690}{[\faFilePdf]}\href{https://github.com/GAIR-NLP/MoPS}{[\faGithub]}}
  {\textit{ACL 2024}}
  {\underline{Yan Ma}, Yu Qiao, Pengfei Liu}
  {}
  \Paper
  {Weak-to-Strong Reasoning \small \href{https://arxiv.org/abs/2407.13647}{[\faFilePdf]}\href{https://github.com/GAIR-NLP/weak-to-strong-reasoning}{[\faGithub]}}
  {\textit{EMNLP Findings 2024}}
  {Yuqing Yang, \underline{Yan Ma}, Pengfei Liu}
  \Paper
  {OlympicArena: Benchmarking Multi-discipline Cognitive Reasoning for Superintelligent AI \href{https://arxiv.org/abs/2406.12753}{[\faFilePdf]}\href{https://github.com/GAIR-NLP/OlympicArena}{[\faGithub]}}
  {\textit{NeurIPS Dataset \& Benchmark Track 2024}}
  {Zhen Huang, Zengzhi Wang, Shijie Xia, ..., \underline{Yan Ma}, ..., Shaoting Zhang, Dahua Lin, Yu Qiao, Pengfei Liu}
  \resumeSubHeadingListEnd

%-----------编程技能-----------------
\section{编程技能}
  \resumeSubHeadingListStart
    \item{
      \textbf{语言}{: Python, Bash, Markdown, Latex}
      \hfill
      \textbf{技术栈}{: Verl, FSDP, VLLM/SGLang, NeoVim, Tmux}
    }
  \resumeSubHeadingListEnd

\end{document}
